\chapter{Rough Volatility}
\label{ch:rough}

\begin{application}[Why This Chapter]
Rough volatility provides an alternative explanation for the slow decay in volatility autocorrelation that HAR (Chapter~\ref{ch:har}) captures with three components and FIGARCH (Chapter~\ref{ch:garch}) captures with fractional integration.
The RFSV model is a parsimonious forecaster competitive with HAR and LSTM.
The Cont--Das counterargument warns against taking roughness as ground truth.
Project 4 (rough vol vs.\ deep learning) directly tests RFSV against universal LSTM.
\end{application}

Chapter~\ref{ch:har} showed that HAR's three-component structure (daily, weekly, monthly) approximates the slow power-law decay of volatility autocorrelations.
Chapter~\ref{ch:garch} showed that FIGARCH (Section~\ref{sec:figarch}) captures the same slow decay with a fractional differencing parameter $d$.
Both approaches are empirically motivated heuristics: they match the shape of the autocorrelation function, but neither explains \emph{why} the decay is so slow.

This chapter presents a different starting point.
Instead of asking ``how should I model the autocorrelation structure?'' it asks: ``what kind of stochastic process would \emph{generate} the autocorrelation structure we observe?''
The answer, proposed by \citet{GJR2018}, is that the log of realized volatility behaves like a \emph{fractional Brownian motion} with a very low Hurst exponent, around $H \approx 0.1$.
This makes the volatility path much rougher (more jagged) than standard Brownian motion ($H = 0.5$).


% ═══════════════════════════════════════════════════
\section{What Is Roughness?}
\label{sec:roughness}
% ═══════════════════════════════════════════════════

The central concept is the \emph{Hurst exponent}, a single number that describes how jagged or smooth a random path looks.

\begin{prereq}[Brownian Motion (Standard)]
A standard Brownian motion $W_t$ is the canonical ``random walk in continuous time.''
Its key properties:
\begin{itemize}[nosep]
  \item $W_0 = 0$.
  \item Increments $W_{t+h} - W_t$ are Gaussian with mean zero and variance $h$.
  \item Increments over non-overlapping intervals are independent.
  \item Paths are continuous but nowhere differentiable (they look jagged at every magnification).
\end{itemize}
If you zoom into a Brownian path, the zoomed-in portion looks statistically identical to the original.
This property is called \emph{self-similarity}, and the scaling factor is $H = 1/2$: rescaling time by a factor $c$ rescales the path by $c^{1/2}$.
\end{prereq}

\begin{prereq}[Fractional Brownian Motion (fBM)]
\emph{Fractional Brownian motion} $B^H_t$, introduced by \citet{MandelbrotVanNess1968}, generalizes standard Brownian motion by allowing the self-similarity exponent $H$ to take any value in $(0, 1)$.
Its key properties:
\begin{itemize}[nosep]
  \item $B^H_0 = 0$, and increments are Gaussian with mean zero.
  \item $\operatorname{Var}(B^H_{t+h} - B^H_t) = h^{2H}$.
  \item The path is self-similar with exponent $H$: rescaling time by $c$ rescales the path by $c^H$.
  \item When $H = 1/2$, fBM reduces to standard Brownian motion.
  \item When $H \neq 1/2$, increments are \emph{not} independent. They are negatively correlated if $H < 1/2$ and positively correlated if $H > 1/2$.
\end{itemize}
The parameter $H$ is called the \emph{Hurst exponent} (after the hydrologist Harold Edwin Hurst, who studied long-range dependence in Nile river levels).
\end{prereq}

The Hurst exponent controls two things simultaneously: the roughness of the path and the correlation structure of increments.

\begin{definition}[Hurst Exponent and Path Regularity]
\label{def:hurst}
For a fractional Brownian motion $B^H_t$ with Hurst exponent $H \in (0,1)$:
\begin{equation}
  \operatorname{Var}(B^H_{t+h} - B^H_t) = h^{2H}
  \label{eq:fbm-variance}
\end{equation}
\begin{itemize}[nosep]
  \item $H$: the Hurst exponent, controlling both path roughness and increment correlation
  \item $h$: the time lag
  \item $h^{2H}$: the variance of the increment over a window of length $h$
  \item When $H < 1/2$: increments are negatively correlated (a move up makes a move down more likely); the path is \emph{rougher} than Brownian motion
  \item When $H = 1/2$: increments are independent; standard Brownian motion
  \item When $H > 1/2$: increments are positively correlated (trending behavior); the path is \emph{smoother} than Brownian motion
\end{itemize}
\end{definition}

\begin{intuition}[Roughness = Anti-Persistence]
Think of $H$ as a dial.
Turn it below $0.5$ and the path becomes ``anti-persistent'': every upward wiggle is likely followed by a downward wiggle, creating a jagged, rapidly oscillating trajectory.
Turn it above $0.5$ and the path becomes persistent: moves tend to continue in the same direction, creating smooth, trend-like trajectories.
At $H = 0.1$, the path is so anti-persistent that it reverses direction almost constantly.
The result is a path so jagged that it looks qualitatively different from standard Brownian motion.
\end{intuition}

Figure~\ref{fig:hurst-paths} shows sample paths at four values of $H$.
This is the most important visual in the chapter.

\begin{figure}[htbp]
\centering
\begin{tikzpicture}
% ── Panel 1: H = 0.1 (very rough) ──
\begin{axis}[
  name=p1,
  width=0.48\textwidth, height=4.5cm,
  title={\small $H = 0.1$ (very rough)},
  title style={at={(0.5,1.0)}, anchor=south},
  xmin=0, xmax=100, ymin=-2.5, ymax=2.5,
  xtick=\empty, ytick={-2,0,2},
  ylabel style={font=\scriptsize},
  tick label style={font=\scriptsize},
  axis lines=left,
  grid=major, grid style={gray!15},
  clip=true,
]
\addplot[warnred, thick, mark=none, smooth] coordinates {
  (0,0.00) (2,0.18) (4,-0.25) (6,0.40) (8,-0.10) (10,0.35)
  (12,-0.30) (14,0.50) (16,-0.15) (18,0.42) (20,-0.38)
  (22,0.55) (24,-0.20) (26,0.30) (28,-0.45) (30,0.60)
  (32,-0.25) (34,0.48) (36,-0.50) (38,0.35) (40,-0.60)
  (42,0.40) (44,-0.35) (46,0.55) (48,-0.48) (50,0.30)
  (52,-0.55) (54,0.65) (56,-0.40) (58,0.50) (60,-0.62)
  (62,0.45) (64,-0.30) (66,0.58) (68,-0.50) (70,0.35)
  (72,-0.45) (74,0.60) (76,-0.55) (78,0.42) (80,-0.38)
  (82,0.50) (84,-0.60) (86,0.38) (88,-0.42) (90,0.55)
  (92,-0.50) (94,0.45) (96,-0.35) (98,0.48) (100,-0.40)
};
\end{axis}

% ── Panel 2: H = 0.3 (rough) ──
\begin{axis}[
  name=p2, at={(p1.east)}, anchor=west, xshift=0.8cm,
  width=0.48\textwidth, height=4.5cm,
  title={\small $H = 0.3$ (rough)},
  title style={at={(0.5,1.0)}, anchor=south},
  xmin=0, xmax=100, ymin=-2.5, ymax=2.5,
  xtick=\empty, ytick={-2,0,2},
  tick label style={font=\scriptsize},
  axis lines=left,
  grid=major, grid style={gray!15},
  clip=true,
]
\addplot[keyorange, thick, mark=none, smooth] coordinates {
  (0,0.00) (4,0.30) (8,0.10) (12,0.55) (16,0.20)
  (20,0.65) (24,0.30) (28,0.70) (32,0.40) (36,0.85)
  (40,0.50) (44,0.25) (48,0.60) (52,0.15) (56,0.50)
  (60,0.80) (64,0.45) (68,0.90) (72,0.55) (76,0.30)
  (80,0.70) (84,0.35) (88,0.75) (92,0.40) (96,0.85)
  (100,0.50)
};
\end{axis}

% ── Panel 3: H = 0.5 (standard BM) ──
\begin{axis}[
  name=p3, at={(p1.below south west)}, anchor=north west, yshift=-0.8cm,
  width=0.48\textwidth, height=4.5cm,
  title={\small $H = 0.5$ (standard Brownian motion)},
  title style={at={(0.5,1.0)}, anchor=south},
  xmin=0, xmax=100, ymin=-2.5, ymax=2.5,
  xtick=\empty, ytick={-2,0,2},
  ylabel style={font=\scriptsize},
  tick label style={font=\scriptsize},
  axis lines=left,
  grid=major, grid style={gray!15},
  clip=true,
]
\addplot[defblue, thick, mark=none, smooth] coordinates {
  (0,0.00) (5,0.25) (10,-0.10) (15,0.40) (20,0.80)
  (25,0.50) (30,0.90) (35,0.60) (40,1.10) (45,0.75)
  (50,1.20) (55,0.85) (60,0.50) (65,0.80) (70,1.30)
  (75,1.00) (80,0.70) (85,1.15) (90,0.90) (95,1.40)
  (100,1.10)
};
\end{axis}

% ── Panel 4: H = 0.7 (smooth) ──
\begin{axis}[
  name=p4, at={(p3.east)}, anchor=west, xshift=0.8cm,
  width=0.48\textwidth, height=4.5cm,
  title={\small $H = 0.7$ (smooth, trending)},
  title style={at={(0.5,1.0)}, anchor=south},
  xmin=0, xmax=100, ymin=-2.5, ymax=2.5,
  xtick=\empty, ytick={-2,0,2},
  tick label style={font=\scriptsize},
  axis lines=left,
  grid=major, grid style={gray!15},
  clip=true,
]
\addplot[intgreen, thick, mark=none, smooth] coordinates {
  (0,0.00) (7,0.20) (14,0.50) (21,0.85) (28,1.15)
  (35,1.40) (42,1.55) (49,1.60) (56,1.50) (63,1.30)
  (70,1.15) (77,1.05) (84,1.10) (91,1.25) (100,1.50)
};
\end{axis}

\end{tikzpicture}
\caption{Sample paths of fractional Brownian motion at four values of the Hurst exponent $H$.
Top left ($H = 0.1$, red): extremely rough, rapidly oscillating.
Top right ($H = 0.3$, orange): rough but with visible short-range structure.
Bottom left ($H = 0.5$, blue): standard Brownian motion, the familiar random walk.
Bottom right ($H = 0.7$, green): smooth, trending trajectory.
Empirically, $\log \RV$ behaves like the top-left panel.}
\label{fig:hurst-paths}
\end{figure}

\begin{keyidea}[Roughness Means $H < 1/2$]
A process is called ``rough'' when its Hurst exponent satisfies $H < 1/2$.
The path is more jagged than Brownian motion, and its increments are negatively correlated.
The lower $H$ is, the rougher the path.
Empirical log-volatility has $H \approx 0.1$, which is far into the rough regime.
\end{keyidea}


% ═══════════════════════════════════════════════════
\section{Volatility Is Rough}
\label{sec:vol-is-rough}
% ═══════════════════════════════════════════════════

With the concept of roughness in hand, the central empirical claim of the rough volatility literature can be stated precisely.

\citet{GJR2018} studied the time series of $\log \RV_t$ (5-minute realized variance, as defined in Chapter~\ref{ch:rv}) across a wide range of assets: equity indices, individual stocks, and foreign exchange.
For each asset, they estimated the Hurst exponent $H$ of the $\log \RV$ series.
The finding was striking.

\begin{keyresult}[\citet{GJR2018}: Volatility Is Rough]
Across equity indices, individual stocks, and FX, the Hurst exponent of $\log \RV_t$ is consistently around $H \approx 0.1$, far below the $H = 0.5$ of standard Brownian motion.
This means that log-volatility paths are much rougher (more jagged, more rapidly oscillating) than any standard diffusion model would predict.
\end{keyresult}

To put $H = 0.1$ in context: standard stochastic volatility models (Heston, SABR) assume that the volatility process is driven by a standard Brownian motion, which has $H = 0.5$.
An $H$ of $0.1$ means the volatility path reverses direction roughly five times more frequently than a standard diffusion would.

\subsection{How to Estimate $H$: The Variogram Method}

The estimation approach in \citet{GJR2018} uses the scaling relationship in Equation~\eqref{eq:fbm-variance}.
If $X_t$ is an fBM with exponent $H$, then the variance of its increments scales as $h^{2H}$.
In log-log coordinates, this becomes a straight line.

\begin{definition}[Variogram Estimator of $H$]
\label{def:variogram}
For a time series $X_1, X_2, \ldots, X_T$, define the empirical $q$-th moment of increments at lag $h$:
\begin{equation}
  m(q, h) = \frac{1}{T - h} \sum_{t=1}^{T-h} |X_{t+h} - X_t|^q
  \label{eq:variogram}
\end{equation}
\begin{itemize}[nosep]
  \item $X_t$: the time series (here, $\log \RV_t$)
  \item $h$: the lag (number of days)
  \item $q$: the moment order (typically $q = 2$ for the variance, or $q = 1$ for the mean absolute increment)
  \item $T$: the number of observations
\end{itemize}
If $X_t$ behaves like fBM with exponent $H$, then $m(q, h) \propto h^{qH}$.
Taking logs: $\log m(q, h) = qH \cdot \log h + \text{const}$.
Regressing $\log m(q, h)$ on $\log h$ across multiple lags gives a slope of $qH$, from which $H$ is recovered.
\end{definition}

\begin{workedexample}{Estimating $H$ from a Variogram}
You have $T = 1{,}000$ daily observations of $\log \RV_t$ for the S\&P~500.
Using $q = 2$, you compute the average squared increment at lags $h = 1, 2, 5, 10, 20$ days:

\medskip
\begin{center}
\begin{tabular}{@{} c c c @{}}
\toprule
\textbf{Lag $h$ (days)} & \textbf{$m(2, h)$} & \textbf{$\log_{10} m(2,h)$} \\
\midrule
1   & 0.050 & $-1.301$ \\
2   & 0.059 & $-1.229$ \\
5   & 0.074 & $-1.131$ \\
10  & 0.088 & $-1.056$ \\
20  & 0.104 & $-0.983$ \\
\bottomrule
\end{tabular}
\end{center}
\medskip

\textbf{Step 1: Set up the regression.}
Plot $\log_{10} m(2,h)$ against $\log_{10} h$:
\begin{center}
\begin{tabular}{@{} c c @{}}
\toprule
$\log_{10} h$ & $\log_{10} m(2,h)$ \\
\midrule
0.000 & $-1.301$ \\
0.301 & $-1.229$ \\
0.699 & $-1.131$ \\
1.000 & $-1.056$ \\
1.301 & $-0.983$ \\
\bottomrule
\end{tabular}
\end{center}

\textbf{Step 2: Fit the slope.}
The OLS slope through these five points is approximately $0.245$.

\textbf{Step 3: Recover $H$.}
Since $\text{slope} = qH = 2H$, we get $H = 0.245 / 2 = 0.12$.

This is close to the $H \approx 0.1$ reported by \citet{GJR2018}.
The slow increase of $m(2,h)$ with lag (i.e., the low slope) reflects the anti-persistent, rapidly oscillating nature of log-volatility.
For comparison, a standard Brownian motion would give a slope of $2 \times 0.5 = 1.0$, meaning the variance of increments grows linearly with lag.
\end{workedexample}

\begin{keyresult}[\citet{BLP2022}: Cross-Asset Universality]
\citet{BLP2022} extend the analysis of \citet{GJR2018} to a broad cross-section of asset classes (equities, equity indices, FX, fixed income, commodities) and confirm that $H \approx 0.1$ is universal.
The Hurst exponent does not vary meaningfully across asset classes, geographies, or time periods.
This universality suggests that $H \approx 0.1$ reflects a deep structural property of volatility dynamics, not a peculiarity of any particular market.
\end{keyresult}


% ═══════════════════════════════════════════════════
\section{The RFSV Forecasting Formula}
\label{sec:rfsv}
% ═══════════════════════════════════════════════════

The observation that $\log \RV$ behaves like fBM with $H \approx 0.1$ has a direct forecasting payoff.
If you know the process is fBM, the optimal linear forecast is known in closed form.
This gives a volatility forecasting model called RFSV (Rough Fractional Stochastic Volatility).

The intuition is simple.
For fBM, the best forecast of the future value is a weighted average of all past observations, where the weights are determined by $H$.
When $H$ is small (rough regime), the weights decay slowly, meaning distant past observations still carry information.
This slow weight decay is what produces the long-memory behavior that HAR approximates with three components.

\begin{prereq}[Conditional Expectation for Gaussian Processes]
If $(X_1, \ldots, X_T, X_{T+1})$ is jointly Gaussian, the conditional expectation $\E[X_{T+1} \mid X_1, \ldots, X_T]$ is a linear combination of the past values.
The weights are determined by the covariance structure of the process.
For fBM, the covariance function is known analytically, so the optimal weights can be computed exactly.
\end{prereq}

\begin{definition}[RFSV Forecast]
\label{def:rfsv}
Model $\log \RV_t$ as fractional Brownian motion with Hurst exponent $H$ and a constant mean $\mu$:
\begin{equation}
  \log \RV_t = \mu + \sigma_v \, B^H_t
  \label{eq:rfsv-model}
\end{equation}
\begin{itemize}[nosep]
  \item $\log \RV_t$: the natural logarithm of day $t$'s realized variance
  \item $\mu$: the unconditional mean of log-volatility
  \item $\sigma_v$: the volatility-of-volatility (a scaling constant)
  \item $B^H_t$: fractional Brownian motion with exponent $H$
\end{itemize}
The optimal one-step-ahead forecast, given observations through day $T$, is:
\begin{equation}
  \widehat{\log \RV}_{T+1} = \sum_{k=0}^{T-1} w_k \, \log \RV_{T-k}
  \label{eq:rfsv-forecast}
\end{equation}
\begin{itemize}[nosep]
  \item $w_k$: the weight on the observation $k$ days in the past
  \item The weights $\{w_k\}$ are determined by the covariance structure of fBM (i.e., by $H$)
  \item They sum to 1 and decay as a power law: $w_k \propto k^{H - 3/2}$ for large $k$
  \item When $H = 0.1$, the decay is $k^{-1.4}$, which is slow enough that observations from weeks ago still contribute meaningfully
\end{itemize}
\end{definition}

\begin{keyidea}[One Parameter Does the Work of Three]
RFSV achieves HAR-level forecasting accuracy with essentially one free parameter ($H$), because the Hurst exponent $H \approx 0.1$ encodes the entire autocorrelation structure.
HAR uses three coefficients ($\beta_d$, $\beta_w$, $\beta_m$) to approximate the same slow decay.
RFSV derives the weights from first principles given $H$.
The fact that both approaches perform similarly is evidence that the slow power-law decay really is the dominant structure in volatility dynamics.
\end{keyidea}

\begin{intuition}[Why RFSV Weights Decay Slowly]
Consider two forecasting extremes.
If $H = 0.5$ (standard BM), increments are independent and only the most recent observation matters; the weights drop to zero immediately.
If $H = 0.01$ (extremely rough), the process is so anti-persistent that every past reversal carries information about the current level; the weights decay very slowly.
At $H = 0.1$, you are close to the second extreme.
The weight on an observation from 22 days ago (one month) is still roughly 15--20\% of the weight on yesterday's observation.
This is why HAR's monthly component ($\RV^{(m)}$) carries a large coefficient in Chapter~\ref{ch:har}: it is picking up the long tail of the RFSV weight function.
\end{intuition}

\begin{workedexample}{Comparing RFSV Weights to HAR Structure}
Consider a simplified RFSV with $H = 0.1$ and 22 past observations.
The weights decay approximately as $w_k \propto k^{-1.4}$.
Normalizing so they sum to 1:

\medskip
\begin{center}
\begin{tabular}{@{} l c c @{}}
\toprule
\textbf{Lag group} & \textbf{Sum of RFSV weights} & \textbf{HAR analog} \\
\midrule
Lag 1 (yesterday)     & $\approx 0.35$ & $\beta_d \approx 0.36$ \\
Lags 2--5 (rest of week)  & $\approx 0.30$ & Embedded in $\beta_w \approx 0.28$ \\
Lags 6--22 (rest of month) & $\approx 0.35$ & Embedded in $\beta_m \approx 0.28$ \\
\bottomrule
\end{tabular}
\end{center}
\medskip

The RFSV weight distribution across lag groups is strikingly similar to the HAR coefficient values from \citet{Corsi2009}.
HAR is an efficient three-bin approximation of the smooth RFSV weight function.
\end{workedexample}


% ═══════════════════════════════════════════════════
\section{Rough Volatility for Pricing}
\label{sec:rough-pricing}
% ═══════════════════════════════════════════════════

This guide focuses on realized volatility estimation and forecasting, not on derivatives pricing.
But the rough volatility framework has had its greatest quantitative-finance impact in the pricing domain, so a brief mention is warranted for completeness.

\begin{prereq}[The Volatility Smile Problem]
Standard stochastic volatility models (e.g., Heston) can generate ``smiles'' in implied volatility across strike prices.
But they struggle to simultaneously fit two empirical patterns: (1) the steep short-maturity smile in S\&P~500 options and (2) the term structure of the VIX implied volatility (the ``VVIX'' smile).
Fitting both at once has been a long-standing challenge in quantitative finance.
\end{prereq}

\citet{BFG2016} introduced the \emph{rough Bergomi} model, the first pricing model built on rough volatility.
In the rough Bergomi model, the variance process is driven by fractional Brownian motion with $H \approx 0.07$, rather than the standard Brownian motion used in models like Heston.

Two key results from the pricing literature:

\begin{enumerate}[nosep]
  \item \textbf{Rough Bergomi} \citep{BFG2016}: the short-maturity behavior of the implied volatility smile is controlled by the Hurst exponent $H$. With $H \approx 0.07$, the model generates steep short-dated smiles that match market data far better than Heston. Simulation is required for pricing (no closed-form solution), but the model is parsimonious.

  \item \textbf{Quadratic rough Heston}: a tractable variant that jointly fits SPX option smiles and VIX option smiles, a combination that no standard model can achieve. The roughness parameter is again $H \approx 0.05$--$0.1$.
\end{enumerate}

\begin{keyidea}[Forecasting vs.\ Pricing: Same $H$, Different Models]
The forecasting model (RFSV) and the pricing models (rough Bergomi, quadratic rough Heston) both use $H \approx 0.1$ but serve different purposes.
RFSV forecasts $\RV_{t+1}$ from past $\RV$ values.
The pricing models compute option prices under risk-neutral dynamics.
This chapter is about forecasting; the pricing models are mentioned here because they provide independent confirmation that the roughness parameter $H \approx 0.1$ is relevant across different uses of volatility modeling.
\end{keyidea}


% ═══════════════════════════════════════════════════
\section{Fact or Artefact?}
\label{sec:fact-or-artefact}
% ═══════════════════════════════════════════════════

The rough volatility paradigm rests on an empirical observation: $\log \RV_t$ has $H \approx 0.1$.
But $\RV_t$ is not the true spot volatility $\sigma^2_t$.
It is a noisy estimate of it, contaminated by the microstructure noise discussed in Chapter~\ref{ch:noise} (Section~\ref{sec:sampling-frequency}).
Could the apparent roughness be an artefact of this noise?

\citet{ContDas2024} argue that the answer is yes, at least in part.

\begin{keyresult}[\citet{ContDas2024}: Roughness as a Noise Artefact]
Observed roughness of realized volatility estimates is partly a microstructure-noise artefact.
When i.i.d.\ noise contaminates the high-frequency returns used to compute $\RV$, the resulting $\RV$ series appears rougher than the true spot volatility process.
Even if the true $\sigma^2_t$ follows a standard semimartingale with $H = 0.5$, the estimated $\log \RV_t$ can exhibit $H \approx 0.1$ due to the noise channel.
\end{keyresult}

The mechanism is intuitive and illustrated in Figure~\ref{fig:noise-roughness}.

\begin{figure}[htbp]
\centering
\begin{tikzpicture}

% ── Panel 1: True spot volatility (smooth) ──
\begin{axis}[
  name=top,
  width=0.85\textwidth, height=4.5cm,
  title={\small True spot volatility $\sigma^2_t$ (smooth, $H = 0.5$)},
  title style={at={(0.5,1.0)}, anchor=south},
  xmin=0, xmax=100, ymin=0.5, ymax=2.5,
  xtick=\empty,
  ylabel={\scriptsize $\sigma^2_t$},
  ylabel style={font=\scriptsize},
  tick label style={font=\scriptsize},
  axis lines=left,
  grid=major, grid style={gray!15},
]
\addplot[intgreen, very thick, mark=none, smooth] coordinates {
  (0,1.0) (5,1.05) (10,1.15) (15,1.25) (20,1.40)
  (25,1.50) (30,1.55) (35,1.50) (40,1.40) (45,1.35)
  (50,1.30) (55,1.35) (60,1.45) (65,1.55) (70,1.60)
  (75,1.55) (80,1.45) (85,1.35) (90,1.30) (95,1.25) (100,1.20)
};
\end{axis}

% ── Panel 2: Estimated RV (noisy, looks rough) ──
\begin{axis}[
  name=bottom, at={(top.below south west)}, anchor=north west, yshift=-0.6cm,
  width=0.85\textwidth, height=4.5cm,
  title={\small Estimated $\RV_t$ (noisy, \emph{looks} rough, $\hat{H} \approx 0.1$)},
  title style={at={(0.5,1.0)}, anchor=south},
  xmin=0, xmax=100, ymin=0.5, ymax=2.5,
  xtick=\empty,
  xlabel={\scriptsize Trading day},
  ylabel={\scriptsize $\RV_t$},
  ylabel style={font=\scriptsize},
  xlabel style={font=\scriptsize},
  tick label style={font=\scriptsize},
  axis lines=left,
  grid=major, grid style={gray!15},
]
% Underlying smooth curve (dashed)
\addplot[intgreen, dashed, thick, mark=none, smooth] coordinates {
  (0,1.0) (5,1.05) (10,1.15) (15,1.25) (20,1.40)
  (25,1.50) (30,1.55) (35,1.50) (40,1.40) (45,1.35)
  (50,1.30) (55,1.35) (60,1.45) (65,1.55) (70,1.60)
  (75,1.55) (80,1.45) (85,1.35) (90,1.30) (95,1.25) (100,1.20)
};
% Noisy RV path
\addplot[warnred, thick, mark=none] coordinates {
  (0,1.05) (2,0.82) (4,1.20) (6,0.90) (8,1.30) (10,1.00)
  (12,1.35) (14,1.10) (16,1.40) (18,1.15) (20,1.55)
  (22,1.30) (24,1.60) (26,1.35) (28,1.55) (30,1.70)
  (32,1.40) (34,1.65) (36,1.45) (38,1.55) (40,1.25)
  (42,1.50) (44,1.20) (46,1.45) (48,1.30) (50,1.15)
  (52,1.40) (54,1.20) (56,1.50) (58,1.30) (60,1.55)
  (62,1.40) (64,1.65) (66,1.45) (68,1.70) (70,1.50)
  (72,1.65) (74,1.45) (76,1.60) (78,1.35) (80,1.50)
  (82,1.30) (84,1.45) (86,1.20) (88,1.40) (90,1.15)
  (92,1.35) (94,1.10) (96,1.30) (98,1.15) (100,1.25)
};
\end{axis}

% Arrow between panels
\draw[-{Stealth[length=3mm]}, thick, gray]
  ([xshift=3cm]top.south) -- ([xshift=3cm]bottom.north)
  node[midway, right, font=\scriptsize, text=gray, align=left] {add estimation\\noise};

\end{tikzpicture}
\caption{How microstructure noise creates apparent roughness.
Top: the true spot volatility $\sigma^2_t$ is a smooth process ($H = 0.5$).
Bottom: the estimated $\RV_t$ (red) adds i.i.d.\ estimation noise around the true level (dashed green).
The noisy $\RV_t$ series looks rough (anti-persistent, rapidly oscillating) even though the underlying process is smooth.
Applying the variogram estimator to the red series yields $\hat{H} \approx 0.1$, but this reflects noise contamination, not true roughness of $\sigma^2_t$.}
\label{fig:noise-roughness}
\end{figure}

The \citet{ContDas2024} argument proceeds in three steps:

\begin{enumerate}[nosep]
  \item \textbf{Noise adds independent errors.} Each day's $\RV_t$ differs from the true $\IVol_t$ by an estimation error $\eta_t$ (Chapter~\ref{ch:rv}, Equation~\ref{eq:rv-clt}). These errors are approximately independent across days.

  \item \textbf{Independent errors create anti-persistence.} If you add i.i.d.\ noise to a smooth signal, the increments of the noisy series are negatively autocorrelated. An unusually high noise draw today will be followed (on average) by a smaller one tomorrow, creating artificial ``reversals'' in the series.

  \item \textbf{Anti-persistence lowers $\hat{H}$.} The variogram estimator interprets these reversals as roughness and produces $\hat{H} \ll 0.5$.
\end{enumerate}

\begin{warning}[Do Not Equate $\hat{H} = 0.1$ with True Roughness]
The observed $H \approx 0.1$ from $\log \RV_t$ data does not establish that the true spot volatility process $\sigma^2_t$ is rough.
The noise channel documented by \citet{ContDas2024} is a plausible alternative explanation.
The truth may lie in between: some genuine roughness in $\sigma^2_t$, amplified by estimation noise in $\RV_t$.
For forecasting, this distinction does not matter much (see Section~\ref{sec:universal-lstm}).
For model calibration and theoretical claims about the nature of volatility, it matters a great deal.
\end{warning}

\begin{intuition}[The Coin-Flip Analogy]
Suppose you flip a fair coin 100 times and record the running total of heads minus tails.
The resulting path is a standard random walk ($H = 0.5$).
Now suppose someone records your totals but makes small random errors (e.g., occasionally miscounting by $\pm 1$).
If you estimate $H$ from their error-contaminated records, you will get $\hat{H} < 0.5$, because the errors introduce artificial reversals.
The path \emph{looks} rougher than it is.
Cont--Das argue that $\RV_t$ is the error-contaminated record and $\sigma^2_t$ is the true total.
\end{intuition}


% ═══════════════════════════════════════════════════
\section{The Universal LSTM Connection}
\label{sec:universal-lstm}
% ═══════════════════════════════════════════════════

\citet{RosenbaumZhang2022} train a single LSTM (Long Short-Term Memory neural network) on hundreds of individual stocks simultaneously.
The trained network, which they call the ``universal LSTM,'' matches the forecasting performance of the RFSV model and the Quadratic Rough Heston parametric forecaster.

\begin{prereq}[LSTM (Brief)]
An LSTM is a type of recurrent neural network designed to learn sequential patterns.
It processes a time series one step at a time, maintaining an internal ``memory cell'' that selectively remembers and forgets past information.
LSTMs are widely used for time series forecasting.
Chapter~\ref{ch:dl-vol} covers LSTMs in detail.
\end{prereq}

This convergence is suggestive.
Two very different approaches (a parametric fBM model with one parameter and a nonparametric neural network with thousands of parameters) arrive at the same forecast.
The natural interpretation is that both are learning the same underlying statistical regularity in $\RV$ data.

\begin{keyidea}[Practical Equivalence Despite Theoretical Ambiguity]
Whether spot volatility is truly rough or the roughness is an artefact of estimation noise (Section~\ref{sec:fact-or-artefact}), both RFSV and universal LSTM exploit the same statistical regularity in realized volatility estimates.
The practical forecasting value is the same either way.
The debate about the nature of the spot process is theoretically important but does not affect your forecast accuracy.
\end{keyidea}

\citet{RosenbaumZhang2022} also document a ``universality'' property: the LSTM trained on a broad cross-section of stocks transfers well to assets it has never seen, including out-of-sample asset classes.
This is consistent with the cross-asset universality of $H \approx 0.1$ documented by \citet{BLP2022}.

\begin{keyresult}[\citet{RosenbaumZhang2022}: Universal LSTM Matches RFSV]
A single LSTM trained on hundreds of stocks achieves the same forecasting performance as the RFSV parametric model.
Both produce forecasts that are competitive with HAR.
The universal LSTM also transfers to out-of-sample assets, suggesting that the learned kernel is asset-independent, just as the Hurst exponent $H \approx 0.1$ is.
\end{keyresult}

The connection to the rest of this guide is direct.
Chapter~\ref{ch:dl-vol} develops LSTM and transformer-based forecasting models in full detail.
When you reach that chapter, the key question will be: does your deep learning model learn something \emph{beyond} the rough-vol kernel, or is it just a complicated way to recover the same $H \approx 0.1$ decay structure?
If the latter, RFSV with one parameter is preferable on parsimony grounds.
If the former, you need to demonstrate the improvement with a proper out-of-sample test (Chapter~\ref{ch:evaluation}).


% ═══════════════════════════════════════════════════
\section*{Summary}
\addcontentsline{toc}{section}{Summary}
% ═══════════════════════════════════════════════════

\begin{itemize}[nosep]
  \item The \textbf{Hurst exponent} $H$ of a fractional Brownian motion controls path roughness. $H = 0.5$ is standard Brownian motion; $H < 0.5$ is rougher (more jagged); $H > 0.5$ is smoother.

  \item \textbf{Fractional Brownian motion} (fBM) generalizes standard BM by allowing correlated increments: negatively correlated when $H < 0.5$ (anti-persistent), positively correlated when $H > 0.5$ (persistent).

  \item \citet{GJR2018} showed that $\log \RV_t$ behaves like fBM with $H \approx 0.1$ across equity indices, individual stocks, and FX. This is the ``volatility is rough'' result.

  \item $H$ can be estimated from data using the \textbf{variogram method}: regress $\log m(q,h)$ on $\log h$ and divide the slope by $q$.

  \item \citet{BLP2022} confirmed the \textbf{cross-asset universality} of $H \approx 0.1$ across equities, FX, fixed income, and commodities.

  \item The \textbf{RFSV model} forecasts $\log \RV_{t+1}$ as a weighted average of past $\log \RV$ values, with weights derived from the fBM covariance structure. It achieves HAR-level accuracy with essentially one parameter ($H$).

  \item RFSV weights decay as $k^{H - 3/2}$. At $H = 0.1$, the decay is slow ($k^{-1.4}$), explaining why HAR's monthly component carries a large coefficient.

  \item \textbf{Rough Bergomi} \citep{BFG2016} and quadratic rough Heston use $H \approx 0.07$--$0.1$ for option pricing. They fit short-maturity smiles and VIX smiles better than standard models. This guide focuses on forecasting, not pricing.

  \item \citet{ContDas2024} argue that observed roughness ($H \approx 0.1$) is \textbf{partly a microstructure-noise artefact}. Estimation noise in $\RV_t$ creates artificial anti-persistence that the variogram interprets as roughness.

  \item The practical implication: $H \approx 0.1$ from $\log \RV_t$ is a useful forecasting input but does not prove that the true spot volatility $\sigma^2_t$ is rough.

  \item \citet{RosenbaumZhang2022} show that a \textbf{universal LSTM} trained on hundreds of stocks matches RFSV and transfers to unseen assets. Both approaches likely learn the same statistical kernel.

  \item Whether spot vol is truly rough or the roughness is a noise artefact, the \textbf{forecasting value is identical}. The distinction matters for pricing-model calibration and theory, not for forecasting accuracy.

  \item RFSV connects to FIGARCH (Chapter~\ref{ch:garch}, Section~\ref{sec:figarch}) and HAR (Chapter~\ref{ch:har}) as three different parameterizations of the same long-memory phenomenon. It connects forward to deep learning (Chapter~\ref{ch:dl-vol}) as a benchmark that any LSTM must beat to justify its complexity.
\end{itemize}


% ═══════════════════════════════════════════════════
\section*{Key Results}
\addcontentsline{toc}{section}{Key Results}
% ═══════════════════════════════════════════════════

\begin{center}
\renewcommand{\arraystretch}{1.3}
\begin{tabular}{@{} p{3.5cm} l p{6.5cm} @{}}
\toprule
\textbf{Result} & \textbf{Source} & \textbf{Finding} \\
\midrule
Volatility is rough
  & \citet{GJR2018}
  & $\log \RV_t$ behaves like fBM with $H \approx 0.1$ across equity
    indices, stocks, and FX; far below $H = 0.5$ of standard models \\[4pt]
Cross-asset universality
  & \citet{BLP2022}
  & $H \approx 0.1$ is universal across asset classes, geographies,
    and time periods \\[4pt]
Rough Bergomi pricing
  & \citet{BFG2016}
  & First rough-vol pricing model; $H \approx 0.07$ generates steep
    short-maturity implied volatility smiles \\[4pt]
Roughness as artefact
  & \citet{ContDas2024}
  & Microstructure noise in $\RV$ estimates creates apparent roughness;
    $\hat{H} \approx 0.1$ may not reflect true spot-vol dynamics \\[4pt]
Universal LSTM
  & \citet{RosenbaumZhang2022}
  & Single LSTM trained on hundreds of stocks matches RFSV;
    both learn the same universal forecasting kernel \\
\bottomrule
\end{tabular}
\end{center}
